\section{Multi-dataset Validation}
\label{sec:multidataset}

The original two-dataset evaluation (SQuAD + PubMedQA) was deliberately
designed to contrast a domain-matched setting against a
domain-mismatched one, but is inadequate for testing whether the
coherence paradox generalizes across question styles, document genres,
and reasoning structures. This section reports an extension to six
datasets covering open-domain factoid QA, multi-hop reasoning, and a
specialized financial domain.

\subsection{Datasets}

\begin{description}
  \item[\textsc{SQuAD}~\citep{rajpurkar2016squad}] reading comprehension
        on Wikipedia paragraphs. Domain-matched single-document QA.
  \item[\textsc{PubMedQA}~\citep{jin2019pubmedqa}] biomedical research
        QA over abstracts. Domain-mismatched single-document QA.
  \item[\textsc{NaturalQuestions}~\citep{kwiatkowski2019nq}] real Google
        queries paired with Wikipedia long-answer paragraphs. Tests
        open-domain factoid retrieval.
  \item[\textsc{TriviaQA}~\citep{joshi2017triviaqa}] trivia questions
        paired with Wikipedia/web evidence. Tests retrieval over
        heterogeneous evidence formats.
  \item[\textsc{HotpotQA}~\citep{yang2018hotpotqa}] multi-hop questions
        requiring reasoning over $\geq$2 paragraphs (distractor split).
        Tests whether the paradox manifests under multi-hop pressure.
  \item[\textsc{FinanceBench}~\citep{islam2023financebench}] open-book
        questions over corporate 10-K filings. Tests a specialized,
        long-document, professional domain.
\end{description}

The complete loader inventory is in \texttt{src/dataset\_loaders.py}.
All loaders share a common \texttt{(documents, qa\_pairs)} interface so
the multi-dataset harness can iterate without per-dataset branching, and
all use a fixed seed (42) for sample selection.

\subsection{Preregistered reporting plan}
\label{sec:multidataset_plan}

For each of the six datasets we will report four quantities, each
averaged over the matched query sample: (i) baseline faithfulness
Faith\textsubscript{base}, (ii) HCPC-v1 faithfulness
Faith\textsubscript{v1}, (iii) HCPC-v2 faithfulness
Faith\textsubscript{v2}, and two derived deltas —
\texttt{paradox\_drop} =
Faith\textsubscript{base} $-$ Faith\textsubscript{v1} (the failure mode)
and \texttt{v2\_recovery} =
Faith\textsubscript{v2} $-$ Faith\textsubscript{v1} (the effectiveness
of selective gating). Per-query NLI scores are written to
\texttt{results/multidataset/summary.csv}; the aggregate table is
generated deterministically from that file.

The paradox is considered to replicate at six-dataset scale if
\texttt{paradox\_drop} is positive on at least four of the six datasets,
with the direction preserved on both the domain-matched (SQuAD,
NaturalQs, TriviaQA) and domain-mismatched (PubMedQA, FinanceBench)
subsets. HCPC-v2's selective gating is considered to help if
\texttt{v2\_recovery} is positive on every dataset where
\texttt{paradox\_drop} is positive. We will report the outcome as
written regardless of direction; a flat or inverted pattern on specific
datasets is informative because it tells us where the coherence account
stops generalizing.

\subsection{Generator stratification (Item~A3)}

The same multi-dataset query set is rerun across three generators
(\texttt{src/generators.py} registry):

\begin{itemize}
  \item \textbf{Mistral-7B-Instruct-v0.2} (Mistral AI). Primary generator.
  \item \textbf{Llama-3-8B-Instruct} (Meta). Cross-architecture replication.
  \item \textbf{Qwen2.5-7B-Instruct} (Alibaba). Tests whether the paradox
        is specific to Western-pretrained model families.
\end{itemize}

The generator-stratified table
(\texttt{results/multidataset/by\_generator.csv}) reports the same three
faithfulness conditions on the multidataset query set, separately for
each generator. The paradox is called \emph{generator-invariant} if the
sign of \texttt{paradox\_drop} is preserved across all three generators
on at least four of the six datasets; a sign flip on any row localizes
the paradox to a specific (dataset, generator) cell rather than to the
method as such. We will report the full cell-level table with confidence
intervals regardless of outcome, and call out any cell where the three
generators disagree on sign.

\subsection{Reproducibility checkpoints}

The multi-dataset harness (\texttt{experiments/run\_multidataset\_validation.py})
is checkpointed at the (dataset, generator, condition) tuple level: each
completed tuple is appended to \texttt{completed\_tuples.json} with its
per-query CSV, and the runner re-enters by skipping completed tuples.
This makes it safe to re-run after partial failure (e.g., Self-RAG
unavailable on M4) without recomputing finished tuples.
